% Gerado por docs/backlog/gera_backlog.py. Não editar à mão.
\begin{table}[htbp]
\centering
\caption{Requisitos novos propostos a partir da RSL 1}
\label{tab:req-novos}
\small
\begin{tabular}{lp{9.2cm}p{3cm}}
\hline
\textbf{ID} & \textbf{Requisito} & \textbf{Origem} \\
\hline
RF11 & Busca e filtros: buscar entidades e evidências do caso por nome, identificador e tipo. & RSL 1, A4 TD14 \\
RF12 & Alerta como pista: todo alerta exige confirmação ou descarte pelo investigador, com justificativa registrada na auditoria. & RSL 1, A4 TD18 e TA3 \\
RF13 & Extração assistida por LLM: entidades e campos extraídos de documentos, sempre com trecho e fonte citados e validação de esquema. & RSL 1, A4 TD13 \\
RF14 & Priorização de alertas: ordenar por severidade, com a política de erro (falso positivo ou falso negativo) declarada por tipologia. & RSL 1, A4 TD11 \\
RF15 & Análise temporal e geográfica: linha do tempo e mapa dos eventos do caso. & RSL 1, A4 TD15 \\
RNF12 & Resiliência da coleta: retentativa, controle de taxa, registro de falha por fonte e detecção de mudança de formato. & RSL 1, A4 TD4; A3 TA3 \\
\hline
\end{tabular}
\fonte{Elaborada pelo autor, a partir da síntese da RSL 1 (Seção~\ref{sec:resultados-rsl1}).}
\end{table}

\begin{table}[htbp]
\centering
\caption{Histórias do \textit{backlog} por prioridade MoSCoW}
\label{tab:moscow-resumo}
\small
\begin{tabular}{lrr}
\hline
\textbf{Prioridade} & \textbf{Histórias} & \textbf{Feitas no protótipo} \\
\hline
Must have & 11 & 11 \\
Should have & 13 & 5 \\
Could have & 24 & 8 \\
Won't have (nesta entrega) & 8 & 1 \\
\hline
\end{tabular}
\fonte{Elaborada pelo autor. MoSCoW de H40 a H57: decisão do autor, a validar pelo orientador.}
\end{table}

\begin{small}
\begin{longtable}{lp{7.9cm}p{1.3cm}cp{1.9cm}}
\caption{\textit{Backlog} do protótipo}\rotulounico{tab:backlog}\\
\hline
\textbf{ID} & \textbf{História} & \textbf{Req.} & \textbf{MoSCoW} & \textbf{Situação} \\ \hline\endhead
\multicolumn{5}{l}{\textbf{E1. Coleta de dados abertos}} \\
H01 & Como investigador, quero coletar empresas e sócios dos dados abertos do CNPJ e contratos do PNCP, para reunir num caso dados oficiais de fontes diferentes. & RF01, RNF10 & M & feito \\
H15 & Como investigador, quero saber, em cada coleta, quais fontes falharam e por quê, com retentativa automática, para não tirar conclusões de dados incompletos. & RNF12 & S & parcial \\
H19 & Como investigador, quero coletar também do Portal da Transparência e do TSE, para cruzar contratos com sanções e doações de campanha. & RF01 & C & feito \\
H25 & Como investigador, quero importar outros dados obtidos por ordem judicial (RIF do COAF, extratos telefônicos), separados do dado aberto, para analisá-los no mesmo caso. & RF10 & W & fora desta entrega \\
H27 & Como investigador, quero coletar de serviços da dark web, para investigar mercados ilícitos. & novo & W & fora desta entrega \\
H50 & Como investigador, quero coletar os bens declarados ao TSE pelos candidatos do caso em cada eleição, com a custódia de cada arquivo, para montar a linha do tempo patrimonial. & RF01, RF05 & S & aguarda orientador \\
H54 & Como investigador, quero importar dados bancários recebidos por ordem judicial no leiaute oficial do Banco Central, com validação de cada arquivo e conferência do período e das contas autorizados, para analisá-los no caso sem digitação. & RF10 & C & a fazer \\
\multicolumn{5}{l}{\textbf{E2. Entidades e vínculos}} \\
H02 & Como investigador, quero que a mesma empresa ou pessoa seja unificada entre as fontes, com sugestão (e não fusão) quando só o nome coincide, para não confundir homônimos. & RF02 & M & feito \\
H03 & Como investigador, quero ver o grafo de vínculos e o caminho entre duas entidades, para identificar sócios ocultos e cadeias societárias. & RF03 & M & feito \\
H41 & Como investigador, quero um alerta quando a mesma pessoa administrar muitas empresas abertas em pouco tempo, para identificar possível laranja profissional. & RF03 & S & a fazer \\
H42 & Como investigador, quero um alerta quando o controle de uma empresa do caso passar por várias pessoas jurídicas ou terminar em sócio estrangeiro, para identificar cadeias societárias em camadas. & RF03 & S & a fazer \\
\multicolumn{5}{l}{\textbf{E3. Alertas e análise assistida}} \\
H04 & Como investigador, quero alertas de tipologia (T1 a T3) que mostrem a regra, o motivo e as evidências com a fonte, para entender por que cada alerta foi gerado. & RF04, RNF05 & M & feito \\
H12 & Como investigador, quero marcar cada alerta como confirmado ou descartado, com justificativa registrada na auditoria, para que o alerta seja tratado como pista e a decisão fique com um humano. & RF12 & S & feito \\
H13 & Como investigador, quero ver os alertas ordenados por severidade, com a política de erro de cada tipologia declarada, para priorizar o que analisar primeiro. & RF14 & S & feito \\
H21 & Como investigador, quero ver os eventos do caso numa linha do tempo e, quando houver endereço, num mapa, para perceber padrões temporais e geográficos. & RF15 & C & a fazer \\
H22 & Como investigador, quero extrair entidades de documentos com apoio de LLM, sempre com o trecho e a fonte citados, para acelerar a leitura sem perder a verificabilidade. & RF13, RNF05 & C & parcial \\
H28 & Como investigador, quero classificar automaticamente textos e imagens por tipo de crime, para triar grandes volumes. & novo & W & fora desta entrega \\
H30 & Como investigador, quero analisar conteúdo em vários idiomas, com gírias e jargão criminal, para não perder informação em fontes estrangeiras. & novo & W & fora desta entrega \\
H31 & Como investigador, quero um alerta quando a empresa contratada não estiver com situação cadastral ativa no CNPJ, para identificar contratação irregular. & RF04 & S & feito \\
H32 & Como investigador, quero um alerta de possível fracionamento de compras (várias contratações do mesmo objeto, com o mesmo fornecedor, abaixo do limite de dispensa), para identificar fuga à licitação. & RF04 & C & feito \\
H33 & Como investigador, quero um alerta quando a empresa contratada estiver no cadastro de empresas sancionadas, para identificar contratação de empresa impedida. & RF04 & C & feito \\
H34 & Como investigador, quero um alerta quando um sócio da empresa contratada for servidor do órgão contratante, para identificar conflito de interesse. & RF04 & C & feito \\
H35 & Como investigador, quero um alerta de termos aditivos acima do limite legal, para identificar alteração indevida de contratos. & RF04 & C & feito \\
H36 & Como investigador, quero que o texto de editais seja analisado em busca de cláusulas restritivas, com o trecho citado, para identificar direcionamento da licitação. & RF13 & C & parcial \\
H39 & Como investigador, quero um alerta quando um doador de campanha (ou empresa de um sócio doador) for contratado pelo ente do candidato eleito, para identificar possível retribuição de doação. & RF04 & C & feito \\
H40 & Como investigador, quero um alerta quando muitas empresas estiverem registradas no mesmo endereço, mesmo sem contrato público, para identificar possível domicílio de fachada. & RF04, RF03 & S & a fazer \\
H43 & Como investigador, quero um alerta quando um MEI, microempresa ou empresa de pequeno porte receber em contratos públicos, no ano, mais que o limite de receita do seu regime, para identificar empresa de fachada. & RF04 & S & a fazer \\
H45 & Como investigador, quero um alerta quando uma empresa recém-aberta contratada tiver sócio que possivelmente é servidor, eleito ou PEP, para identificar desvio por empresa criada para contratar. & RF04 & C & a fazer \\
H46 & Como investigador, quero um alerta quando um eleito ou PEP for possivelmente sócio de empresa contratada pelo próprio ente durante o mandato, para identificar autocontratação. & RF04 & C & aguarda orientador \\
H47 & Como investigador, quero um alerta quando a contratada compartilhar sócio-administrador ou endereço com empresa sancionada, ou tiver sócio expulso da administração, para identificar empresa sucessora de sancionada. & RF04 & C & a fazer \\
H48 & Como investigador, quero um alerta de combinação quando a mesma entidade acumular vários alertas, com a lista de cada um e das evidências, para priorizar o que tem mais sinais juntos. & RF14, RF12 & C & a fazer \\
H49 & Como investigador, quero sinais fracos (contato compartilhado, sócio em faixa etária atípica, excesso de CNAEs) que só pesem no alerta de combinação, para não gerar alertas isolados de baixo valor. & RF04 & C & a fazer \\
H51 & Como investigador, quero ver a evolução dos bens declarados de um candidato entre eleições, e a renda pública conhecida quando houver, com alerta só quando a variação superar a renda, para identificar indício de enriquecimento ilícito. & RF05 & S & aguarda orientador \\
H52 & Como investigador, quero um alerta quando o candidato pagar, com recursos de campanha, empresa de atividade sem relação com campanha, recém-aberta ou com sinais de fachada, para identificar desvio de recursos de campanha. & RF04 & C & aguarda orientador \\
H53 & Como investigador, quero um alerta quando um servidor doar a campanha valor desproporcional à sua remuneração pública, para identificar possível doação por interposta pessoa. & RF04 & W & aguarda orientador \\
H56 & Como investigador, quero saber se o sócio de uma contratada recebe benefício social, para identificar possível laranja. & RF03 & W & fora desta entrega \\
H57 & Como investigador, quero um alerta quando uma entidade sem fins lucrativos que recebe convênio for recém-aberta, estiver no CEPIM ou tiver dirigente que possivelmente é eleito, PEP, servidor do concedente ou sócio de empresas do mesmo ramo, para identificar desvio por convênio. & RF04, RF03 & C & a fazer \\
\multicolumn{5}{l}{\textbf{E4. Prova e cadeia de custódia}} \\
H05 & Como investigador, quero que cada carga coletada seja guardada com hash SHA-256, carimbo de tempo UTC e origem, e que eu possa verificar a integridade depois, para que a prova se sustente em juízo. & RF07 & M & feito \\
H08 & Como investigador, quero exportar o caso (entidades, alertas, evidências e custódia) em JSON e Markdown, para anexar ao relatório de investigação. & RF09 & M & feito \\
H14 & Como perito, quero que a captura registre também a versão do conector e que eu possa recoletar a mesma fonte e comparar, para reconstituir como a prova foi obtida. & RF07 & S & feito \\
H24 & Como analista, quero exportar o caso num formato padronizado de troca (por exemplo, CASE/UCO), para levar os dados a outras ferramentas. & RNF09 & C & a fazer \\
H38 & Como analista, quero exportar o grafo do caso para i2 Analyst's Notebook e Maltego, para usar as ferramentas que o órgão já tem. & RNF09 & C & parcial \\
\multicolumn{5}{l}{\textbf{E5. Casos e busca}} \\
H07 & Como investigador, quero criar casos e fazer toda consulta dentro de um caso, com justificativa, para organizar a investigação e limitar a finalidade. & RF08 & M & feito \\
H09 & Como investigador, quero buscar entidades do caso por nome, CNPJ ou CPF parcial e filtrar por tipo, para encontrar rapidamente o que interessa. & RF11 & M & feito \\
H37 & Como investigador, quero anotar entidades e alertas e dividir tarefas do caso com a equipe, para trabalhar em conjunto na mesma investigação. & RF08 & C & feito \\
\multicolumn{5}{l}{\textbf{E6. Segurança, auditoria e conformidade}} \\
H06 & Como auditor, quero uma trilha só de acréscimo, encadeada por hash, de quem consultou o quê, quando e com qual justificativa, para detectar uso indevido. & RNF02 & M & feito \\
H10 & Como responsável pelo sistema, quero que todo acesso seja feito por usuário identificado e com perfil, para controlar quem usa a ferramenta. & RNF01 & M & feito \\
H11 & Como delegado responsável, quero que a ferramenta só faça coleta passiva de dados abertos, sem interagir com alvos, para não configurar infiltração virtual sem ordem judicial. & RNF04 & M & feito \\
H16 & Como desenvolvedor, quero análise estática de segurança e os testes como barreira no pipeline de integração, para que código gerado por IA não entre sem verificação. & RNF11 & S & parcial \\
H17 & Como delegado responsável, quero definir a finalidade e o prazo de retenção de cada caso, para aplicar a minimização dos dados coletados. & RNF04 & S & feito \\
H23 & Como responsável pelo sistema, quero MFA, acesso por caso e criptografia em repouso, para proteger dados de investigação. & RNF01 & C & parcial \\
H26 & Como investigador, quero coletar por infraestrutura não atribuível à polícia, para não alertar o alvo. & RNF03 & W & fora desta entrega \\
H55 & Como delegado responsável, quero que o dado sob ordem judicial fique numa base separada, com chave própria, marcado com o número da decisão, fora das exportações por padrão e com cada leitura auditada, para não misturá-lo ao dado aberto. & RF10, RNF02, RNF04 & C & a fazer \\
\multicolumn{5}{l}{\textbf{E7. Monitoramento e evolução}} \\
H18 & Como investigador, quero ser avisado quando mudar o quadro societário de uma empresa do caso ou surgir contrato novo, para acompanhar alvos sem refazer a coleta. & RF06 & C & feito \\
H29 & Como investigador, quero consultas de vínculo em segundos sobre grafos com milhões de nós, para trabalhar com a base completa do CNPJ. & RNF06 & W & feito \\
H44 & Como investigador, quero um alerta quando nenhum fundador continuar na empresa ou o quadro societário mudar pouco antes de um contrato, para identificar troca de laranjas. & RF03, RF06 & C & a fazer \\
\hline
\multicolumn{5}{p{14cm}}{\footnotesize Fonte: elaborada pelo autor. A evidência de cada história está em \texttt{docs/backlog/backlog.md}. MoSCoW de H40 a H57: decisão do autor, a validar pelo orientador.}\\
\end{longtable}
\end{small}
